\documentclass{report}
\usepackage{amsmath,amssymb}
\setlength{\parindent}{0mm}
\setlength{\parskip}{1em}
\begin{document}
\begin{center}
\rule{6in}{1pt} \
{\large Kees Vuik \\
{\bf On the Convergence of Shifted Laplace Preconditioner Combined with Multigrid Deflation}}

Delft University of Technology \\ Delft Institute of Applied Mathematics \\ Mekelweg 4 \\ 2628 CD Delft \\ The Netherlands
\\
{\tt c.vuik@tudelft.nl}\\
Abdul Sheikh\\
Domenico Lahaye\end{center}

In this paper, a solver for the discretized time-harmonic wave equation
in heterogeneous media
is considered.
The underlying equation governs wave propagations and scattering
phenomena arising in many area's, e.g., in aeronautics, geophysics, and optical
problems. In particular, we look for efficient iterative
solution techniques for the Helmholtz equation
discretized by finite difference discretizations. Since the number of
grid points per wavelength should be sufficiently large for accurate
solutions, for very high wavenumbers the discrete problem becomes
extremely large, thus prohibiting the use of direct solvers.
However, since the coefficient matrix is sparse,
iterative solvers are an interesting alternative.

Starting point of our research is the multilevel Krylov method proposed in [1]
which consists of a combination of the Shifted Laplace Preconditioner (SLP) given in [2]
and a second level preconditioner based on Deflation [3]. It is well known that
the coefficient matrix of the discretized Helmholtz equation has eigenvalues both
in the positive and negative complex plane. Such a spectrum leads to a very slow
convergence of Krylov type solvers. Combining the matrix with a Shifted
Laplace Preconditioner
leads to a spectrum which is enclosed in a circle in the positive complex plane.
If no damping is included this circle passes through the origin.
It appears that the resulting
method is scalable in step size $h$ but the number of iterations increases linearly
with the increase of the wave number. The reason for this increase is that the number
of eigenvalues close to zero increase considerably for large wave numbers. Motivated
by the success of multigrid, where a smoother and coarse grid solver are combined,
Erlangga and Nabben [1] use a combination of SLP and projection Deflation
in a multilevel way.
From their paper it appears that the number of iterations is much smaller
and furthermore the method seems to be scalable also in the wave number. To gain insight
in this method, we analyze the method proposed in more detail.

In this paper we perform a Rigorous Fourier Analysis of a two-level variant of the
multilevel Krylov method for the Helmholtz equation proposed in [1]. The
distinct feature of the solver analyzed is the two-grid deflation of the Shifted
Laplace Preconditioner at each step of an outer Krylov subspace acceleration.
Our analysis of 1D and 2D models reveals two properties of the solver. The
first is that with deflation, the solver performs better than without deflation.
However, above a certain large wave number, the solver
depends again linearly on the wave number. The reason is that the spectrum
contains more small eigenvalues with increasing wave number. The second
is that the Shifted Laplace Preconditioner can be made arbitrarily diagonally
dominant by increasing the imaginary part of the shift without paying any
penalty in the number of Krylov iterations.
These properties are verified by numerical experiments on Helmholtz problems with
heterogeneous media and may lead to faster methods to approximate the
Shifted Laplace Preconditioner.

\section*{References}
[1] Y.A. Erlangga and R. Nabben. On a multilevel Krylov method for the Helmholtz equation
preconditioned by shifted Laplacian. Electronic Transactions on Numerical
Analysis (ETNA),
31:403�424, 2008.
\newline
[2] Y.A. Erlangga, C.W. Oosterlee, and C. Vuik. A novel multigrid based
preconditioner for
heterogeneous Helmholtz problems. SIAM J. Sci. Comput, 27:1471�1492, 2006.
\newline
[3] J.M. Tang, S.P. MacLachlan, R. Nabben, and C. Vuik
A Comparison of Two-Level Preconditioners Based on Multigrid and Deflation
SIAM. J. Matrix Anal. and Appl., 31: 1715-1739, 2010


\end{document}
